\begin{ouverture}[grandeurs]{12}{Les aires : unités et formules}{Partie 3 · Grandeurs et mesures}
{Comparaisons · cm², dm², m² · Carré et rectangle}
\objectif{Comparer et mesurer des aires par pavage.}
\objectif{Connaître et convertir les unités d’aire usuelles.}
\objectif{Calculer l’aire d’un carré et d’un rectangle.}
\prerequis{Aires par comptage (CM1) et multiplication des décimaux.}
\end{ouverture}
\pagestyle{otMain}\essentiel\fontsize{11}{15}\selectfont
\phantomsection\label{cours-12-1}
\essSection{1}{Des carreaux aux unités officielles}
\begin{definition}Le centimètre carré ($\mathrm{cm}^2$) est l’aire d’un carré de 1 cm de côté. Le décimètre carré ($\mathrm{dm}^2$) et le mètre carré ($\mathrm m^2$) suivent la même idée : carrés de 1 dm et de 1 m de côté.\end{definition}
\label{fig-cours-12-1-cours-dm2}\figureGeo{cours-dm2}

\begin{propriete}Dans un carré de 1 dm de côté, on range 10 rangées de 10 petits carrés de 1 cm : $1\ \mathrm{dm}^2=100\ \mathrm{cm}^2$. De même, $1\ \mathrm m^2=100\ \mathrm{dm}^2$. D’une unité d’aire à la suivante, le rapport est 100, pas 10.\end{propriete}
\begin{savoirfaire}[Convertir des aires]Convertir : $3\ \mathrm{dm}^2=300\ \mathrm{cm}^2$ ; $4{,}5\ \mathrm m^2=450\ \mathrm{dm}^2$ ; $200\ \mathrm{cm}^2=2\ \mathrm{dm}^2$. Citer le rapport 100 à chaque étape.\end{savoirfaire}
\clearpage
\phantomsection\label{cours-12-2}
\essSection{2}{Les formules du carré et du rectangle}
\begin{propriete}L’aire d’un rectangle est longueur $\times$ largeur, mesurées dans la même unité : un rectangle de 7 cm sur 4 cm pave $7\times4=28$ carreaux de $1\ \mathrm{cm}^2$. L’aire du carré est côté $\times$ côté.\end{propriete}
\begin{exemple}Un carré de 6 cm de côté : aire $6\times6=36\ \mathrm{cm}^2$ et périmètre $4\times6=24$ cm. Les deux grandeurs ne se confondent pas : l’une pave, l’autre entoure.\end{exemple}
\begin{savoirfaire}[Découper ou compléter]Pour une figure en L, découper en rectangles : un L de côtés 8 et 6 avec une encoche de 5 sur 4 se traite par différence : $8\times6-5\times4=28$ unités d’aire, ou par découpage en deux rectangles.\end{savoirfaire}
\begin{attention}Les longueurs doivent partager la même unité avant la formule : 2 m sur 50 cm, c’est $200\times50=10\,000\ \mathrm{cm}^2$, soit $1\ \mathrm m^2$, et non « $2\times50=100$ ».\end{attention}

\clearpage\rubrique{Découverte}
\begin{blocLivre}{ex-12-1}\begin{exercice}[1]{Compter puis calculer}
Un rectangle quadrillé compte 4 rangées de 7 carreaux de $1\ \mathrm{cm}^2$. Donner son aire par comptage, puis par la formule.
\label{fig-ex-12-1}\figureGeo{aire-1}
\end{exercice}\end{blocLivre}
\begin{blocLivre}{ex-12-2}\begin{exercice}[1]{Carré et rectangle}
Calculer l’aire d’un carré de 6 cm de côté, puis d’un rectangle de 9 cm sur $4{,}5$ cm.
\end{exercice}\end{blocLivre}
\begin{blocLivre}{ex-12-3}\begin{exercice}[1]{Aire ou périmètre ?}
Pour un carré de 6 cm de côté, calculer le périmètre et l’aire, puis préciser l’unité de chacun.
\end{exercice}\end{blocLivre}
\begin{blocLivre}{ex-12-4}\begin{exercice}[1]{Les conversions}
Convertir : $3\ \mathrm{dm}^2$ en $\mathrm{cm}^2$ ; $4{,}5\ \mathrm m^2$ en $\mathrm{dm}^2$ ; $200\ \mathrm{cm}^2$ en $\mathrm{dm}^2$.
\end{exercice}\end{blocLivre}
\begin{blocLivre}{ex-12-5}\begin{exercice}[1]{Pourquoi 100 ?}
Expliquer pourquoi $1\ \mathrm{dm}^2=100\ \mathrm{cm}^2$ alors que $1$ dm $=10$ cm.
\label{fig-ex-12-5}\figureGeo{aire-5}
\end{exercice}\end{blocLivre}

\rubrique{Entraînement}
\begin{blocLivre}{ex-12-6}\begin{exercice}[2]{Retrouver un côté}
L’aire d’un rectangle vaut $54\ \mathrm{cm}^2$ et sa longueur 9 cm. Quelle est sa largeur ? Et le côté d’un carré de $49\ \mathrm{cm}^2$ ?
\end{exercice}\end{blocLivre}
\begin{blocLivre}{ex-12-7}\begin{exercice}[2]{La figure en L}
Une pièce en L s’obtient en retirant un rectangle de 5 m sur 4 m d’un rectangle de 8 m sur 6 m. Calculer son aire de deux façons.
\label{fig-ex-12-7}\figureGeo{aire-7}
\end{exercice}\end{blocLivre}
\begin{blocLivre}{ex-12-8}\begin{exercice}[2]{Les unités d’abord}
Calculer l’aire d’un rectangle de 2 m sur 50 cm, en $\mathrm{cm}^2$ puis en $\mathrm m^2$.
\end{exercice}\end{blocLivre}
\begin{blocLivre}{ex-12-9}\begin{exercice}[2]{Même périmètre, aires rangées}
Trois rectangles ont un périmètre de 20 cm : $9\times1$, $7\times3$ et $5\times5$. Vérifier, puis les ranger par aire croissante.
\end{exercice}\end{blocLivre}
\begin{blocLivre}{ex-12-10}\begin{exercice}[2]{Le carré agrandi}
On triple le côté d’un carré de 4 cm. Par combien l’aire est-elle multipliée ?
\end{exercice}\end{blocLivre}

\rubrique{Problèmes}
\begin{blocLivre}{ex-12-11}\begin{exercice}[3]{La chambre à moquette}
Une chambre rectangulaire mesure $4{,}2$ m sur $3{,}5$ m. La moquette coûte 18 euros le mètre carré. Calculer l’aire de la chambre puis le prix de la moquette.
\end{exercice}\end{blocLivre}
\begin{blocLivre}{ex-12-12}\begin{exercice}[3]{Le potager et l’allée}
Un terrain carré de 12 m de côté contient une allée rectangulaire de 12 m sur 2 m ; le reste est cultivé. Quelle surface est cultivée ?
\end{exercice}\end{blocLivre}
\begin{blocLivre}{ex-12-13}\begin{exercice}[3]{Les dalles de la terrasse}
Une terrasse rectangulaire mesure 6 m sur 4 m. On la pave avec des dalles carrées de 50 cm de côté. Combien de dalles faut-il ?
\end{exercice}\end{blocLivre}
\begin{blocLivre}{ex-12-14}\begin{exercice}[3]{Deux annonces}
Deux appartements affichent : A « $45\ \mathrm m^2$ » et B « un rectangle de 9 m sur 5,2 m ». Lequel est le plus grand, et de combien ?
\end{exercice}\end{blocLivre}
\begin{blocLivre}{ex-12-15}\begin{exercice}[3]{Vrai ou faux ?}
« Si deux rectangles ont la même aire, ils ont le même périmètre. » Trancher avec les rectangles $12\times3$ et $9\times4$.
\end{exercice}\end{blocLivre}

\clearpage
